\section{Background}

\subsection{Analog design automation}

Traditional analog circuit automation has long depended on expert knowledge and manual optimization. Early design automation tools relied on predefined topology libraries selected by human designers, focusing mainly on parameter sizing rather than structural innovation. Subsequent AI-based methods, including reinforcement learning~\cite{zhao2020deep,9576505} and graph neural networks~\cite{liu2021parasitic,dong2023cktgnn}, expanded automation to limited topology generation and device tuning. Despite these advances, most works are limited to simplified device models, constrained search spaces, and expert-crafted heuristics. Automating topology synthesis remains challenging \cite{10937153} due to the vast design space, nonlinear performance trade-offs, and the nuanced judgment that human designers intuitively apply.

\subsection{Context management in AI agent}

Effective use of LLMs in complex design workflows requires careful management of limited context windows through isolation, compression, and reflection\cite{mei2025survey}. Isolated contexts allow agents to focus on specific subtasks without interference, compressed contexts summarize essential details to conserve space, and reflective contexts capture insights or mistakes from prior reasoning to improve future performance~\cite{chang2025sagallm}. Techniques such as summarization, pruning, and retrieval-based re-injection help maintain coherence across multi-step reasoning while preventing information overload~\cite{packer2023memgpt}. Through dynamic distillation and reuse of relevant information, the agent sustains continuity, recalls prior design choices, and applies analogical reasoning to new problems~\cite{lu2025scaling}.

Recent works such as ReAct~\cite{yao2022react} and Reflexion~\cite{shinn2023reflexion} demonstrate concrete implementations of reflective context mechanisms, combining action traces with self-evaluation to enhance long-horizon reasoning. Their iterative feedback processes highlight the importance of structured context updates for robust, multi-step decision making.



\subsection{Agentic frameworks in circuit design}

Recent research has explored combining LLMs with circuit simulators for analog design automation. However, many existing approaches still rely on prompt engineering~\cite{vungarala2024spicepilot} or rule-based scripts~\cite{zhang2025analogxpert} without true feedback loops or specification-driven reasoning. Their validation is often conducted using idealized device models~\cite{chen2025menter} rather than practical PDK environments, limiting real-world applicability. Furthermore, context length and knowledge retention are underutilized, preventing effective cross-task learning and simulation-grounded feedback~\cite{lai2025analogcoder,chen2024artisan}. 

In addition, most prior systems lack rigorous quantitative evaluation, do not test generality across multiple specification-driven tasks, and rarely release complete code or experimental pipelines, making independent reproduction and fair comparison difficult.

This work addresses these limitations through a self-evolving, simulation-grounded framework that integrates multi-context fusion, introspective feedback, and numerical optimization, while providing open-source code.